\section{Comparative evaluation}
\subsection{The four different text representations}
The classical models share Logistic Regression, the data split and validation. Vocabulary sizes and feature geometry still differ because the representations work differently: for example, a single word produces several character $n$-grams. DistilBERT \cite{sanh2019distilbert} has a different architecture, so it serves as a separate reference.

\begin{center}
\begin{minipage}{\linewidth}
\centering\small
\captionof{table}{Model input representations.}
\begin{tabularx}{\linewidth}{@{}L{0.27\linewidth}Y@{}}
\toprule
\tabhead Model & Input representation \\
\midrule
Word TF--IDF & Word unigrams and bigrams. \\
Character TF--IDF & Character 3--5-grams. \\
BPE TF--IDF & \texttt{o200k\_base} token-ID unigrams and bigrams. \\
DistilBERT & Uncased WordPiece; maximum 128 tokens. \\
\bottomrule
\end{tabularx}
\end{minipage}
\end{center}

For all three TF--IDF models, I chose \(C\in\{0.1,1,10\}\) by validation Macro-F1; each selected \(C=10\) on both datasets. DistilBERT used three epochs, a learning rate of \(2\times10^{-5}\), effective batch size 16, weight decay 0.01 and warmup ratio 0.10.

\subsection{What changes in an attacked input?}
Each attack samples eligible positions reproducibly, then the clean and attacked examples gets stored together with the edits made.

\begin{center}
\begin{minipage}{\linewidth}
\centering\small
\captionof{table}{Attack families and where they apply.}
\begin{tabularx}{\linewidth}{@{}L{0.29\linewidth}Y@{}}
\toprule
\tabhead Family & Transformation and eligibility \\
\midrule
Homoglyph & Replace mapped letters with curated Greek or Cyrillic lookalikes. \\
Invisible & Insert a default-ignorable control after an alphanumeric character. \\
Compatibility & Use alternate forms whose NFKC normalization recovers the source. \\
Composition & Decompose a precomposed character when it has a canonical decomposition. \\
Mixed & At each eligible position, sample one applicable transformation. \\
\bottomrule
\end{tabularx}
\end{minipage}
\end{center}

The attack strengths tested are 5\%, 10\%, 20\%, 40\% and 75\% of eligible positions.

\subsection{Coverage of the methods}
A nominal budget is meaningful only if the attack can effectively change the text samples. The four broad families reach every test example at 20\% (I'll take this as the main strength from now on); canonical composition reaches none (Table~\ref{tab:coverage}). Its unchanged test score therefore provides no evidence of robustness, because no test input has an eligible canonical decomposition. Naturally accented text could give a different result.

\begin{center}
\begin{minipage}{\linewidth}
\centering\small\setlength{\tabcolsep}{3pt}
\captionof{table}{Coverage and average edits per example at 20\% strength. B77 denotes BANKING77.}
\begin{tabular}{lrrrr}
\toprule
Attack & B77 cov. & B77 changes & TREC cov. & TREC changes \\
\midrule
homoglyph & 1.00 & 4.86 & 1.00 & 3.26 \\
invisible & 1.00 & 8.54 & 1.00 & 5.64 \\
compatibility & 1.00 & 8.85 & 1.00 & 5.88 \\
composition & 0.00 & 0.00 & 0.00 & 0.00 \\
mixed & 1.00 & 8.85 & 1.00 & 5.88 \\
\bottomrule
\end{tabular}

\label{tab:coverage}
\end{minipage}
\end{center}

Budgets are relative to each family's eligible positions so the same percentage does not produce the same number of edits for different attack kinds. 

\subsection{Five test-time defenses}
The models are evaluated without retraining. Each wrapper processes both clean and attacked inputs, making any clean-text harm visible. NFC and NFKC follow Unicode normalization rules; neither is a general solution to cross-script lookalikes \cite{unicode15,unicode39}.

\begin{table*}[!tbp]
\centering\small
\caption{Defense scope. ``Full repair'' names the most complete wrapper in this experiment.}
\begin{tabularx}{\linewidth}{@{}L{1.35in}Y Y@{}}
\toprule
\tabhead Wrapper & Operation & Intended coverage \\
\midrule
NFC & Canonical composition. & Canonical equivalence; untested here because composition coverage is zero. \\
NFKC & Compatibility normalization. & The generated compatibility substitutions. \\
Control removal & Remove default-ignorable characters. & The selected invisible insertions. \\
NFKC + removal & Apply both operations in sequence. & Compatibility and invisible changes. \\
Full repair & Add restricted non-ASCII-to-ASCII confusable repair. & Also the curated homoglyphs; assumes English/ASCII targets. \\
\bottomrule
\end{tabularx}
\end{table*}


\subsection{How the result measures work}
\label{sec:measures}
Let \(N\) be the number of test examples and \(K\) the number of classes. For class \(k\), \(TP_k\) counts correct predictions of that class, \(FP_k\) counts incorrect predictions of it, and \(FN_k\) counts its missed examples.


\minisection{Accuracy} Correct predictions divided by \(N\); each example has equal weight.

\minisection{Macro-F1} Precision and recall are balanced within each class, then all classes have equal weight:
\[
M=\frac{1}{K}\sum_{k=1}^{K}\frac{2TP_k}{2TP_k+FP_k+FN_k}.
\]
From now on, \(M\) denotes Macro-F1. A class with no correct predictions receives F1 zero. Accuracy weights frequent classes more heavily, while Macro-F1 prevents them from dominating the average. That's why it is useful in these multiclass classification tasks.

\minisection{Macro-F1 loss} \(M_{\mathrm{raw,clean}}-M_{\mathrm{raw,attacked}}\). A loss of 0.10 is ten F1 percentage points, not a 10\% relative loss. 

\minisection{Attack success rate (ASR)} Newly incorrect examples divided by examples correct when clean. It measures how often an attack overturns an initially correct prediction; baseline errors are excluded. ASR is undefined if no examples are correct when clean. For example: if 80 of 100 examples are correct when clean and an attack overturns 20 of those 80 predictions, ASR is \(20/80=25\%\). It is not the attacked error rate: the 20 original errors are outside its denominator, and some of those could become correct after attack (which would be by sheer luck and it's not worth investigating in the current study).


\minisection{Clean-score delta} \(M_{\mathrm{defended,clean}}-M_{\mathrm{raw,clean}}\). A negative value means that the wrapper harms clean performance.



\subsection{Defense recovery \texorpdfstring{\(R\)}{R}}
\label{sec:recovery-definition}
Recovery measures the fraction of the attack-induced Macro-F1 loss that a wrapper restores:
\begin{equation}
R=\frac{M_{\mathrm{defended,attacked}}-M_{\mathrm{raw,attacked}}}
        {M_{\mathrm{raw,clean}}-M_{\mathrm{raw,attacked}}}.
\label{eq:recovery}
\end{equation}

If \(R=0\), the attacked score is unchanged; \(R=0.5\) means half of the lost Macro-F1 is recovered; and \(R=1\) means the clean baseline is reached. Values are not clipped so repair can make the score worse (\(R<0\)) or raise it above the clean baseline (\(R>1\)), lastly Recovery is undefined when the raw attack causes no positive loss.

For each condition, I calculate the metrics on the matched test examples, but to keep tidyness and avoid entire pages of only tables for some of them I averaged the scores of the different attacks families. The overall attacked score gives equal weight to homoglyph, invisible, compatibility and mixed attacks. Composition is excluded from that average because it changes no inputs.
